% Refactor twin claim proof file for row claim_3_2 -- "AcyclicIffTopologicalOrder".
% Created by the write_tex_proof worker for the claim_3_2_no_finite refactor.
%
% This twin will replace the original tex proof at Phase 7 cleanup. The
% LN's iterative parent-free-node construction (preserved verbatim in
% `claim_3_2_proof_AcyclicIffTopologicalOrder.tex`) presupposes
% finiteness of $J \cup V$ -- it explicitly cites "since $G_i$ is acyclic
% ... and finite, it has a node $v_i$ with $\Pa^{G_i}(v_i) = \emptyset$"
% -- so it cannot serve as the proof of the no-finiteness variant of the
% Lem. The twin below proves the same equivalence by the Szpilrajn
% order-extension route, which does not enumerate vertices and so works
% for arbitrary (possibly infinite) $J \cup V$.
%
% Compiles standalone via the `subfiles` package.

\documentclass[main]{subfiles}

\begin{document}
% ref: claim_3_2 (proof twin, with statement restated for readability)
% title: AcyclicIffTopologicalOrder (refactor: no-finiteness variant)
\def\rowref{claim\_3\_2}
\phantomsection\label{claim_3_2}

% --- statement (restated) ---------------------------------------------------
\begin{Lem}[AcyclicIffTopologicalOrder]
        A CDMG  $G=(J,V,E,L)$ is acyclic if and only if it has a topological order.
\end{Lem}

% --- proof ------------------------------------------------------------------
\begin{proof}
This is the \emph{refactor twin}. The original proof, preserved in
\texttt{claim\_3\_2\_proof\_AcyclicIffTopologicalOrder.tex}, constructs
a topological order by iteratively picking a parent-free node from the
induced subgraph on the remaining vertices -- a step which explicitly
invokes finiteness of $J \cup V$. The statement of the Lem makes no
finiteness assumption, so we re-prove the $(\Rightarrow)$ direction via
\emph{Szpilrajn's order-extension theorem} applied to the
``reachable-by-a-directed-walk'' preorder, a route which never enumerates
vertices and therefore works for arbitrary $J \cup V$. Both proofs
establish the same equivalence; the choice is a route change driven by
the removal of the finiteness hypothesis.

\medskip
\noindent\emph{$(\Leftarrow)$ Existence of a topological order implies acyclicity.}
Suppose $<$ is a topological order of $G$ (\refrow{def_3_8}): a total
order on $J \cup V$ satisfying $v \in \Pa^G(w) \implies v < w$. If there
were a non-trivial directed walk $v = v_0 \tuh v_1 \tuh \cdots \tuh v_n = v$
with $n \geq 1$ and $v \in J \cup V$, then $v_{i-1} \in \Pa^G(v_i)$ for
each $i$, so the parent-lt clause gives $v_0 < v_1 < \cdots < v_n$, which
collapses by transitivity to $v_0 < v_n = v_0$, contradicting
irreflexivity of $<$. Hence $G$ has no such walk, i.e.\ $G$ is acyclic
(\refrow{def_3_6}).

\medskip
\noindent\emph{$(\Rightarrow)$ Acyclicity implies the existence of a topological order.}
Assume $G$ is acyclic. We construct a topological order on $J \cup V$
without enumerating its elements.

Define a binary relation $r_0$ on $J \cup V$ by
\[
  r_0(v, w) \;\iff\; \text{there exists a directed walk } v = u_0 \tuh u_1 \tuh \cdots \tuh u_n = w \text{ in } G,\; n \geq 0.
\]
We verify that $r_0$ is a partial order.

\begin{itemize}
\item \emph{Reflexivity.} The trivial walk of length zero witnesses $r_0(v, v)$ for every $v$.

\item \emph{Transitivity.} Given $r_0(x, y)$ and $r_0(y, z)$, concatenate the witnessing walks at their common endpoint $y$. Concatenation preserves directedness (every step of either walk is a forward edge), so the result is a directed walk from $x$ to $z$, witnessing $r_0(x, z)$.

\item \emph{Antisymmetry.} Suppose $r_0(x, y)$ and $r_0(y, x)$, witnessed by directed walks $\pi_1 \colon x \to y$ and $\pi_2 \colon y \to x$. Assume towards a contradiction that $x \neq y$. Then $\pi_1$ has at least one step, so it begins with a forward edge $x \tuh u_1 \in E$; because $E \ins (J \cup V) \times V$ (\refrow{def_3_1}), the source $x$ of this edge lies in $J \cup V$. Now the concatenation $\pi_1$ followed by $\pi_2$ is a directed walk from $x$ to $x$ of length $\geq 1$, i.e.\ a non-trivial directed cycle through $x \in J \cup V$, contradicting acyclicity (\refrow{def_3_6}). Therefore $x = y$.
\end{itemize}

By \emph{Szpilrajn's order-extension theorem} (Szpilrajn, 1930; formalised
in Mathlib as \texttt{extend\_partialOrder}), every partial order extends
to a linear (total) order on the same underlying set. Applied to $r_0$,
this yields a linear order $s$ on $J \cup V$ such that
$r_0(v, w) \implies s(v, w)$ for all $v, w \in J \cup V$.

Take $<$ to be the strict part of $s$:
\[
  v < w \;\iff\; s(v, w) \wedge v \neq w.
\]
We verify the four clauses required for $<$ to be a topological order of
$G$ (\refrow{def_3_8}):

\begin{itemize}
\item \emph{Irreflexivity.} $v < v$ would require $v \neq v$, which is impossible.

\item \emph{Transitivity.} Suppose $v < w$ and $w < u$, i.e.\ $s(v, w)$ with $v \neq w$, and $s(w, u)$ with $w \neq u$. Transitivity of $s$ gives $s(v, u)$. To see $v \neq u$: if $v = u$ then $s(v, w)$ and $s(w, v)$ together force $v = w$ by antisymmetry of the linear order $s$, contradicting $v \neq w$. Hence $v < u$.

\item \emph{Trichotomy.} Given $v, w \in J \cup V$, either $v = w$, or $v \neq w$ and then totality of $s$ supplies $s(v, w)$ or $s(w, v)$, yielding $v < w$ or $w < v$ respectively.

\item \emph{Parent-lt.} Suppose $v \in \Pa^G(w)$, i.e.\ $v \tuh w \in G$. The single forward edge $v \tuh w \in E$ is a directed walk of length one, so $r_0(v, w)$, hence $s(v, w)$. To see $v \neq w$: if $v = w$ then $v \tuh v$ is a non-trivial directed walk from $v$ to $v$, and the source $v$ of this edge lies in $J \cup V$ by $E \ins (J \cup V) \times V$ (\refrow{def_3_1}), contradicting acyclicity (\refrow{def_3_6}). Hence $v < w$.
\end{itemize}

Thus $<$ is a topological order of $G$ in the sense of \refrow{def_3_8},
so $G$ has a topological order.
\end{proof}

\end{document}
